\section{Four Sessions for One Dead Service}
\label{sec:ch14-four-sessions-for-one-dead-service}

\index{partial failure!when the field is not nullable|(}%
Chapter~\ref{ch:second-third-subgraph} stopped the Speakers service and the
graph kept working. Every title came back and every speaker was null, under an
error naming the subgraph and the path it failed at. I called that the blast
radius of a dead subgraph and said it was the set of fields that route to it,
provided those fields are nullable.

\index{errors!the response that chapter carries is the data half}%
One correction to that page before this chapter builds on it. The response
printed there is the \code{data} half and the first of its errors, and there
are five: the fetch failure, and then one for each of the four sessions, because
\code{Speaker.name} is \code{String!} and the router had no speaker to put a
name on. All four of those stop at \code{speaker}, which is nullable, so the
damage is what that chapter said it was. The count is not, and
section~\ref{sec:ch14-how-far-a-null-climbs} needs it.

Stop the Ratings service instead and ask for a field that is not.

\begin{minted}{text}
POST /graphql HTTP/1.1
Host: localhost:3002
Content-Type: application/json

{"query":"{ sessions(first: 10) { nodes { title ratingCount } } }"}
\end{minted}

\begin{minted}[fontsize=\footnotesize,breakanywhere]{json}
{"errors":[{"message":"Failed to fetch from Subgraph 'ratings' at Path 'sessions.nodes'."},
{"message":"Cannot return null for non-nullable field 'Query.sessions.nodes.ratingCount'.",
"path":["sessions","nodes",0,"ratingCount"]}],
"data":{"sessions":{"nodes":null}}}
\end{minted}

\index{blast radius!the whole page for one field}%
The schedule is gone. Not the rating counts, the schedule: four sessions, their
titles, their start times, everything a client would need to render the
program, none of which the Ratings service has ever had anything to do with.
The Sessions service answered correctly and its answer is not in the response.

\index{statement count!work done and discarded}%
It did the work, too. Chapter~\ref{ch:data-without-n-plus-1}'s interceptor
prints one statement in the Sessions service while that request runs. The rows
were read out of SQLite, the titles were in memory, and then a field belonging
to a different service could not be null.

\subsection{The Same Outage, One Exclamation Mark Apart}

\index{nullability!the field that absorbs}%
Ask the same dead service for \code{averageScore} instead:

\begin{minted}{text}
{"query":"{ sessions(first: 10) { nodes { title averageScore } } }"}
\end{minted}

\begin{minted}[fontsize=\footnotesize,breakanywhere]{json}
{"errors":[{"message":"Failed to fetch from Subgraph 'ratings' at Path 'sessions.nodes'."}],
"data":{"sessions":{"nodes":[
{"title":"Schemas That Outlive Their Authors","averageScore":null},
{"title":"Reading a Query Plan","averageScore":null},
{"title":"Paging Without Offsets","averageScore":null},
{"title":"The Bank That Split Its Graph","averageScore":null}]}}}
\end{minted}

\index{partial failure!the good case and the bad one differ by one exclamation mark}%
Same service, same outage, same fetch failure, same one statement in Sessions.
The whole difference is that \code{averageScore} is \code{Float} and
\code{ratingCount} is \code{Int!}, and chapter~\ref{ch:second-third-subgraph}
declared both of them in one chapter, neither as a reliability decision. One of them is a client that renders the program with
blank scores. The other is a client with nothing to render.

\index{errors!the second one is the propagation}%
Count the errors as well. The good case has one, and it is the fetch. The bad
case has two, and the second is not a second failure: it is the consequence of
the first, reported at the position that could not hold it.

\subsection{Which Fields Are Loaded}

\index{non-null!the fields this graph promises}%
Three non-null fields in this graph cross a seam, and they cross it in two
different shapes. Ratings contributes \code{ratingCount} and
\code{feedbackUrl} to a \code{Session} it does not own. Search declares
\code{session} on a document it does own and cannot fill it without asking
Sessions. Ask for \code{feedbackUrl} with Ratings stopped and the page empties
exactly as it did above, with the path naming \code{feedbackUrl} instead.

\index{@requires@\code{@requires}!two ways to fail}%
\code{feedbackUrl} is the worst of the three, because
chapter~\ref{ch:second-third-subgraph} put it behind \code{@requires}. The
router has to collect a title from Sessions before Ratings can be asked for the
url at all, so the field has two services to depend on and one position to fail
at.

\index{blast radius!the shape of the collateral}%
None of this is a router defect and none of it is a Hot Chocolate defect. It is
the specification working exactly as chapter~\ref{ch:schema-design} described,
across a seam that did not exist when those annotations were written. What
changed is not the rule. It is how far the damage travels, and who pays for it.
